% Changed: restructure for a single evidence-bound service-object argument.
% Protected: measured values, protocol distinctions and failure records.
% Meaning-Check: PRESERVED
% Risk-Flags: none
\section{Discussion and Limitations}
\label{sec:discussion}

\subsection{What the Service Contract Establishes}
The reusable object couples layout, evidence, identity, and publication; its benefit depends on matching requests and resident blocks. Token-path certification supplies relative behavioral evidence under the declared greedy execution protocol. It is not an oracle for answer correctness or invariance over arbitrary future inputs. A zero-F1 witness can be preserved exactly, which is why Fig.~\ref{fig:behavior} reports reference utility separately from output-ID agreement.

Query-qualified execution assumes a repeated authorized directory and typed metadata-only observations. It does not cover arbitrary evolving histories or unrestricted new tool evidence. If the number of unique query scopes grows with the request count, qualification has little opportunity to amortize. Production query-repeat frequencies were not measured. Deployment-qualified registration has a different witness origin and reuse unit; the implementation does not automatically switch between two independent endpoint certificate systems.

\subsection{Trust, Authorization, and Timing}
The operator manifest is authoritative for segment stability and sharing permission. A hash establishes content identity, not whether content was correctly declared public. In nine CPU registration cases, seven invalid declarations are rejected; a legitimate declaration and a trusted operator's semantically incorrect declaration are accepted. This helper remains a separate validation test, rather than an integrated live endpoint. The authenticated deployment path and in-engine PSTP controls are evaluated at their stated boundaries.

Authorization-bound namespaces do not remove timing inference within an allowed sharing domain~\cite{iknow2025,auditing2025}. ContractKV does not implement a timing-channel defense. P4 therefore describes cache-state authorization, not a comprehensive application confidentiality boundary.

\subsection{Workload and Platform Boundaries}
The main study uses six independent content sources, two-request concurrency, greedy decoding, and 24-GB GPUs. Tenants and process repeats do not replace new source templates. The preserved-input adapter uses answer-blind optional-context cropping and the declared legacy code scorer. Its Python-output check excludes some legal next-line fragments. The artifact retains omitted spans, inherited L0 routes, and all score changes.

The native-cache probes diagnose three known sources with synchronous instrumentation; their timings are not performance measurements, and their traces cannot retrospectively attribute every historical miss. The separate query-qualified fixed-arrival extension completed five of 54 planned cells and is not used as a capacity result. A vLLM 0.25.1 audit checks relevant integration objects and layout opportunities, not a completed port. Distributed cache domains would require additional lifecycle and publication coordination.

The resulting operating lesson is specific: invest in behavioral qualification where authorized reuse can repay it, measure the cost on the service clock, and account for native reclamation when interpreting a qualified prefix. The original gate remains the principal measured configuration because the tested cost adaptation traded away useful coverage. These conditions define where the reported benefits hold.
